\documentclass[11pt]{article}

\usepackage[margin=1in]{geometry}
\usepackage{amsmath, amssymb}
\usepackage{xcolor}
\usepackage{listings}
\usepackage{hyperref}
\usepackage{array}
\usepackage{booktabs}

\hypersetup{
    colorlinks=true,
    linkcolor=blue!60!black,
    urlcolor=blue!60!black,
    citecolor=blue!60!black
}

\lstdefinestyle{bugpython}{
  language=Python,
  basicstyle=\ttfamily\small,
  keywordstyle=\color{blue!60!black},
  stringstyle=\color{green!40!black},
  commentstyle=\color{gray!70!black},
  showstringspaces=false,
  breaklines=true,
  frame=single,
  rulecolor=\color{gray!30},
  columns=fullflexible,
  keepspaces=true
}

\title{SymPy Bug Report}
\author{}
\date{}
\renewcommand{\bfdefault}{m}

\begin{document}

\maketitle

\section{Bug 20: Conditional Uniform Density Keeps the Wrong Support}

\subsection{Minimal Reproducer}

\medskip

\begin{lstlisting}[style=bugpython]
from sympy import S, integrate, oo, symbols
from sympy.stats import P, Uniform, density, given

x = symbols("x")
U = Uniform("U", 0, 1)
Y = given(U, U > S.Half)
d = density(Y)(x)

print("density:", d)
print("total mass:", integrate(d, (x, -oo, oo)))
print("mass below 1/2 from density:", integrate(d, (x, 0, S.Half)))
print("P(Y < 1/2):", P(Y < S.Half))
\end{lstlisting}

\subsection{Observed Behavior}

\medskip

\begin{lstlisting}[style=bugpython]
density: 2*Piecewise((1, (x >= 0) & (x <= 1)), (0, True))
total mass: 2
mass below 1/2 from density: 1
P(Y < 1/2): 0
\end{lstlisting}

SymPy applies the factor \(2 = 1/P(U > 1/2)\), but the returned density is
still supported on the original interval \([0,1]\).  As a result, the expression
returned by \texttt{density} is not a probability density: it integrates to
\(2\) and assigns mass \(1\) to \([0,1/2]\), even though the conditional
probability query gives \(P(Y < 1/2)=0\).

\subsection{Expected Behavior}

For \(U \sim \operatorname{Uniform}(0,1)\) and \(Y = U \mid U > 1/2\), the
conditional density should be
\[
f_Y(x) =
\begin{cases}
2, & 1/2 < x \le 1,\\
0, & \text{otherwise}.
\end{cases}
\]
Equivalently, the returned density should have total mass \(1\), zero mass below
the conditioning threshold, and value \(0\) at any probe point \(x<1/2\).

\subsection{Mathematical Explanation}

The conditional density is obtained by restricting the original density to the
conditioning event and then renormalizing:
\[
f_{U \mid U>1/2}(x)
  = \frac{f_U(x)\,\mathrm{1}_{\{x>1/2\}}}{P(U>1/2)}.
\]
Since \(f_U(x)=1\) on \([0,1]\) and \(P(U>1/2)=1/2\), this gives density \(2\)
only on the restricted support \((1/2,1]\).

SymPy's result instead represents \(2\mathrm{1}_{[0,1]}(x)\).  This is not an
equivalent form of the conditional density: it is positive on points that
violate the condition, and
\[
\int_{-\infty}^{\infty} 2\mathrm{1}_{[0,1]}(x)\,dx = 2
\qquad\text{while}\qquad
\int_{-\infty}^{\infty} f_{U \mid U>1/2}(x)\,dx = 1.
\]
The discrepancy is therefore a support error and a normalization error, not a
harmless simplification or alternate representation.

\subsection{Source-Level Diagnosis}

The diagnosis identifies the relevant implementation in
\nolinkurl{sympy/stats/crv.py}.

In that file, the relevant methods are
\nolinkurl{ContinuousPSpace.conditional_space}, around lines 446--462, and
\nolinkurl{ContinuousPSpace.compute_density}, around lines 333--340.
The traced call path is
\[
\begin{array}{c}
\texttt{sympy.stats.rv.given} \\
\rightarrow\ \texttt{ContinuousPSpace.conditional\_space} \\
\rightarrow\ \texttt{Density.doit}/\texttt{density} \\
\rightarrow\ \texttt{ContinuousPSpace.compute\_density}.
\end{array}
\]
For the reproducer, the conditional probability space has a
\texttt{ConditionalContinuousDomain} with set
\(\operatorname{Interval.Lopen}(1/2,1)\), so the domain restriction itself is
computed.  The stored density, however, remains the original uniform density
divided only by the conditioning normalizer.

\begin{lstlisting}[style=bugpython]
domain = ConditionalContinuousDomain(self.domain, condition)
norm = domain.compute_expectation(self.pdf, **kwargs)
pdf = self.pdf / norm.xreplace(replacement)
density = Lambda(tuple(domain.symbols), pdf)
return ContinuousPSpace(domain, density)
\end{lstlisting}

This rule renormalizes the original density but does not multiply it by an
indicator for the conditional domain.  The second part of the failure occurs
when \texttt{density(Y)} requests the density of the same variable.  In that
case, \texttt{compute\_density} has no random symbols to marginalize, and
\texttt{ConditionalContinuousDomain.compute\_expectation} returns the stored
expression unchanged when the variable list is empty.  The condition
\(U>1/2\) is therefore never applied to the returned density, producing
\(2\mathrm{1}_{[0,1]}\) instead of \(2\mathrm{1}_{(1/2,1]}\).

The diagnosis suggests that a plausible fix is to ensure that conditional
continuous densities include the condition's support, either by incorporating an
indicator or \texttt{Piecewise} restriction in
\texttt{conditional\_space}, or by making density extraction for a
\texttt{ConditionalContinuousDomain} apply the domain restriction even when no
variables are marginalized.  A general fix would need to avoid double-applying
the original distribution support already present in many distribution PDFs.

\subsection{Additional Failing Instances}

The same failure occurs for a family of conditional uniform densities
\(Y_t = U \mid U>t\) with \(0<t<1\).  For every such threshold, the correct
density is \(1/(1-t)\) only on \((t,1]\) and \(0\) below \(t\).  The verification
cases check both global normalization and local support by integrating the
returned density over \([0,t]\) and evaluating it at \(t/2\).

\begin{center}
\scriptsize
\renewcommand{\arraystretch}{1.3}
\begin{tabular}{@{}>{\raggedright\arraybackslash}p{0.44\linewidth}>{\raggedright\arraybackslash}p{0.23\linewidth}>{\raggedright\arraybackslash}p{0.23\linewidth}@{}}
\toprule
\textbf{Input / Expression} & \textbf{SymPy behavior} & \textbf{Expected behavior} \\
\midrule
\(Y=\operatorname{given}(U,U>1/2)\), \(\operatorname{density}(Y)(x)\)
& \(2\mathrm{1}_{[0,1]}(x)\); total mass \(2\); mass on \([0,1/2]\) is \(1\)
& \(2\mathrm{1}_{(1/2,1]}(x)\); total mass \(1\); zero mass below \(1/2\) \\
\(Y=\operatorname{given}(U,U>1/3)\), \(\operatorname{density}(Y)(x)\)
& \(\frac{3}{2}\mathrm{1}_{[0,1]}(x)\); positive at \(x=1/6\)
& \(\frac{3}{2}\mathrm{1}_{(1/3,1]}(x)\); zero at \(x=1/6\) \\
\(Y=\operatorname{given}(U,U>1/4)\), \(\operatorname{density}(Y)(x)\)
& \(\frac{4}{3}\mathrm{1}_{[0,1]}(x)\); positive at \(x=1/8\)
& \(\frac{4}{3}\mathrm{1}_{(1/4,1]}(x)\); zero at \(x=1/8\) \\
\(Y=\operatorname{given}(U,U>2/3)\), \(\operatorname{density}(Y)(x)\)
& \(3\mathrm{1}_{[0,1]}(x)\); positive at \(x=1/3\)
& \(3\mathrm{1}_{(2/3,1]}(x)\); zero at \(x=1/3\) \\
\(Y=\operatorname{given}(U,U>3/4)\), \(\operatorname{density}(Y)(x)\)
& \(4\mathrm{1}_{[0,1]}(x)\); positive at \(x=3/8\)
& \(4\mathrm{1}_{(3/4,1]}(x)\); zero at \(x=3/8\) \\
\(Y=\operatorname{given}(U,U>1/5)\), \(\operatorname{density}(Y)(x)\)
& \(\frac{5}{4}\mathrm{1}_{[0,1]}(x)\); positive at \(x=1/10\)
& \(\frac{5}{4}\mathrm{1}_{(1/5,1]}(x)\); zero at \(x=1/10\) \\
\bottomrule
\end{tabular}
\end{center}

These are all instances of the same defect: the normalizing factor changes with
\(t\), but the returned support remains the unconditional support \([0,1]\)
instead of the conditional support \((t,1]\).

\subsection{Related Issues Assessment}

The best-effort deduplication check found documentation and background material
for \texttt{density}, \texttt{given}, and conditional probability distributions,
but did not identify a public SymPy issue, pull request, release note, mailing
list thread, or Stack Overflow post for this exact reproducer or the internal
\texttt{ConditionalContinuousDomain} support-retention failure.  The verdict is
therefore a new failure mode with no public precedent.  The bug remains
individually actionable because it has a concise reproducer, a concrete invalid
density, a localized source diagnosis, and a focused regression test.

\subsection{Proposed SymPy Regression Test}

\medskip

\begin{lstlisting}[style=bugpython]
import pytest

from sympy import S, integrate, oo, symbols
from sympy.stats import P, Uniform, density, given


@pytest.mark.parametrize("t", [S.Half, S(1)/3, S(1)/4, S(2)/3, S(3)/4, S(1)/5])
def test_conditional_uniform_density_support_is_restricted(t):
    x = symbols("x")
    U = Uniform("U", 0, 1)
    Y = given(U, U > t)
    d = density(Y)(x)

    assert integrate(d, (x, -oo, oo)) == 1
    assert integrate(d, (x, 0, t)) == 0
    assert d.subs(x, t/2) == 0
    assert P(Y < t) == 0
\end{lstlisting}

\end{document}
